\begin{table}[!htbp]\caption{Economic Value of Mobility-Based Forecasts}\label{tab:economicvalue}\fontsize{11}{12}\selectfont\noindent\begin{tabular}{@{}L{190pt}C{68pt}C{68pt}C{68pt}C{68pt}@{}}\toprule
Strategy & $h=1$ & $h=3$ & $h=6$ & $h=12$\\\midrule
\multicolumn{5}{l}{\textit{Panel A: Annual CER Gain (percentage points)}}\\*
EE & 0.00 & 0.00 & 0.00 & 0.00\\
Historical average & 0.00 & 0.00 & 0.00 & 0.00\\
Buy-and-hold & 0.00 & 0.00 & 0.00 & 0.00\\
\midrule\multicolumn{5}{l}{\textit{Panel B: Annual Sharpe Ratio}}\\*
EE & 0.00 & 0.00 & 0.00 & 0.00\\
Historical average & 0.00 & 0.00 & 0.00 & 0.00\\
Buy-and-hold & 0.00 & 0.00 & 0.00 & 0.00\\
\bottomrule\end{tabular}\par\vspace{7pt}\noindent\begin{minipage}[t]{\linewidth}\exhibitnotestyle\textit{Notes:} This table reports annualized certainty-equivalent return (CER) gains relative to the historical-average strategy (HA) and annualized Sharpe ratios for the EE-based allocation strategy and its HA and buy-and-hold benchmarks. For portfolio allocation, the predictive regression directly forecasts the compounded simple excess return over the $h$-month holding period. The mean--variance investor has risk aversion $\gamma=3$, with equity weights constrained to $[-0.5,1.5]$. Portfolio risk is estimated from the most recent 120 completed, overlapping same-horizon returns. The forecast horizon, holding period, and rebalancing frequency coincide and are given by $h$. HA uses the same risk estimate and portfolio constraints, while buy-and-hold corresponds to an equity weight of one. CERs and Sharpe ratios are annualized by $12/h$ and $\sqrt{12/h}$, respectively. CER gains are percentage points per year, and baseline portfolio returns are before transaction costs.\par\end{minipage}\end{table}
